% Part: lambda-calculus
% Chapter: church-rosser
% Section: beta-reduction

\documentclass[../../../include/open-logic-section]{subfiles}

\begin{document}

\olfileid[tr]{lam}{cr}{b}

\olsection{$\beta$-indirgeme}

\begin{lem}\ollabel{lem:one-par}
  $M \bredone M'$ ise $M \bredpar M'$.
\end{lem}
\begin{proof} $M \bredone M'$ ise bazı $x,N,Q$ için
    $M=(\lambd[x][N])Q$ ve $M'=\Subst{N}{Q}{x}$ olur.
    \olref[pb]{thm:refl} gereği $N\bredpar N$ ve $Q\bredpar Q$ olduğundan,
    \olref[pb]{defn:bredpar}\olref[pb]{defn:bredpar4} gereği doğrudan
    $(\lambd[x][N])Q\bredpar\Subst{N}{Q}{x}$ elde edilir.
\end{proof}

\begin{lem}\ollabel{lem:par-red}
  $M \bredpar M'$ ise $M \bred M'$.
\end{lem}

\begin{proof} $M \bredpar M'$ !!{derivation} üzerinde tümevarımla.
  \begin{enumerate}
  \item Son kural \olref[pb]{defn:bredpar1} ise $M$ ile $M'$ aynı $x$
    değişkenidir ve $x\bred x$.
  \item Son kural \olref[pb]{defn:bredpar2} ise bazı $x,N,N'$ için
    $M=\lambd[x][N]$, $M'=\lambd[x][N']$ ve $N\bredpar N'$ olur.
    Tümevarım varsayımına göre $N\bred N'$. Dolayısıyla $N\bred N'$
    indirgemesindeki aynı $\bredone$ daraltmaları kullanılarak
    $\lambd[x][N]\bred\lambd[x][N']$ elde edilir.
  \item Son kural \olref[pb]{defn:bredpar3} ise bazı $P,Q,P',Q'$ için
    $M=PQ$, $M'=P'Q'$, $P\bredpar P'$ ve $Q\bredpar Q'$ olur. Tümevarım
    varsayımına göre $P\bred P'$ ve $Q\bred Q'$. Önce $P\bred P'$, sonra
    $Q\bred Q'$ indirgemesi uygulanarak $PQ\bred P'Q'$ elde edilir.
  \item Son kural \olref[pb]{defn:bredpar4} ise bazı $x,N,N',Q,Q'$ için
    $M=(\lambd[x][N])Q$, $M'=\Subst{N'}{Q'}{x}$,
    $N\bredpar N'$ ve $Q\bredpar Q'$ olur. Tümevarım varsayımına göre
    $Q\bred Q'$ ve $N\bred N'$. Önce $N\bred N'$, sonra $Q\bred Q'$ ve
    son olarak $(\lambd[x][N'])Q'$ ifadesinin
    $\Subst{N'}{Q'}{x}$'e daraltılmasıyla
    $(\lambd[x][N])Q\bred\Subst{N'}{Q'}{x}$ elde edilir.
  \end{enumerate}
\end{proof}

\begin{lem}\ollabel{lem:str}
  $\bred$, $\bredpar$ bağıntısını içeren en küçük geçişli bağıntıdır.
\end{lem}

\begin{proof}
  $\xred$, $\bredpar$ bağıntısını içeren en küçük geçişli bağıntı olsun.

  $\bred \subseteq \xred$: $M\bred M'$, yani
  $M\ident M_1\bredone\dots\bredone M_k\ident M'$ olduğunu varsayın.
  \olref{lem:one-par} gereği
  $M\ident M_1\bredpar\dots\bredpar M_k\ident M'$. $\xred$, $\bredpar$
  bağıntısını içerip geçişli olduğundan $M\xred M'$.

  $\xred \subseteq \bred$: $M\xred M'$, yani
  $M\ident M_1\bredpar\dots\bredpar M_k\ident M'$ olduğunu varsayın.
  \olref{lem:par-red} gereği
  $M\ident M_1\bred\dots\bred M_k\ident M'$. $\bred$ geçişli olduğundan
  $M\bred M'$.
\end{proof}

\begin{thm}\ollabel{thm:cr}
  $\bred$ Church--Rosser özelliğini sağlar.
\end{thm}

\begin{proof}
  Doğrudan \olref[dap]{thm:str}, \olref[pb]{thm:cr} ve \olref{lem:str}
  sonuçlarından çıkar.
\end{proof}

\end{document}
